\scanpage{551}
\mspage{\unsure{281}}

bijection canonique
\[
\tag{33}
\widetilde I \xrightarrow{\sim} V(I)^*
\overset{\unsure{\mathrm{df}}}{=} V(I)\setminus\{0\}
\]
en associant à une partition $I=I'\cup I''$, avec
$I',I''$ de cardinal $2$, la \inserted{seule} permutation
\deleted{\inserted{(paire)}} de $I$
\deleted{formée de} qui \unsure{invarie} $I',I''$ et \unsure{induise} sur
chacun \unsure{sa} permutation non triviale. (Les éléments
de $\mathfrak{S}_I$ obtenus ainsi sont aussi les \uline{permutations}
\uline{paires} \uline{d'ordre $2$}, \deleted{\unsure{)}} qui avec l'identité forment
donc le sous-groupe $V(I)$).

On trouve, si $i\in I$, d'où $I\setminus\{i\}$ de
cardinal $3$, \deleted{\unsure{que}} une application canonique
\[
\tag{34}
I\setminus\{i\}\longrightarrow\widetilde I
\]
en associant à $j\in I\setminus\{i\}$ la partition
de $I$ en $\{i,j\}$ et en son complémentaire.
Cette bijection étant compatible avec le
transport des structures, on en conclut
que l'isom. correspondant
\[
\tag{35}
\mathfrak{S}_{I\setminus\{i\}}\simeq\mathfrak{S}_{\widetilde I}
\]
\deleted{\unsure{qui}} est \deleted{\unsure{compatible}} aussi le composé
\[
\tag{36}
\underset{\substack{\text{stabilisateur de}\\ i\text{ dans }\mathfrak{S}_I}}
{\mathfrak{S}_{I\setminus\{i\}}}
\hookrightarrow\mathfrak{S}_I\longrightarrow\mathfrak{S}_{\widetilde I},
\]
donc les quatre sous-groupes symétriques
d'indice $3$ $\mathfrak{S}_{I\setminus\{i\}}$ de $\mathfrak{S}_I$ $(i\in I)$\unsure{.|,}

\scanpage{552}
sont canoniquement \unsure{isomorphes} à $\mathfrak{S}_{\widetilde I}$
--- donc entre eux, \deleted{\unsure{mais}} par des isom. qui
\unsure{$\overline{\mathrm m}$} sont \deleted{\unsure{pas|par}} d'ailleurs \deleted{\illegible{1}} déduits des bijections
canoniques \deleted{\unsure{égales}}
\[
  (\deleted{\illegible{1}}\inserted{37})\qquad
  I\setminus\{i\}\xrightarrow{\sim}I\setminus\{j\}
  \qquad(i\ne j,\ i,j\in I)
\]
canoniques, égales à \deleted{l'identité}
\inserted{la \uline{transposition} (non à l'identité)} sur $I\setminus\{i,j\}$.
\inserted{Cette bijection en effet est le composé}
\deleted{(En effet, ces bijections \unsure{ne forment pas un}}
\[
  I\setminus\{i\}\xrightarrow{\sim}\widetilde I
  \xrightarrow{\sim}I\setminus\{j\}
  \qquad\text{---}
\]
\deleted{\unsure{syst.} transitif d'isom.)}
et pour $i,j$ variables, ces isom. forment un
syst. transitif d'isom.

Revenons à l'action de $\mathfrak{S}_4$ sur
$M_{\unsure{0},4}^{!}$, la clef de la \unsure{compréhension}
est donnée par ceci :

\uline{Proposition} Considérons le \unsure{s-groupe}
$V(\deleted{\unsure{4}}\inserted{I})$ de $\mathfrak{S}_{\deleted{\unsure{4}}\inserted{I}}$ formé de $1$ et des involutions
paires. Alors l'action de $V(\deleted{\unsure{4}}\inserted{I})$ sur $M_{0,\deleted{\unsure{4}}\inserted{I}}^{!}$
est triviale, \unsure{\&|Si} donc \deleted{$V(I)$} $\mathfrak{S}_I$ agit
sur $M_{0,I}^{!}$ par l'intermédiaire de
$\mathfrak{S}_I/V(I)=\mathfrak{S}_{\widetilde I}$.
\figure{[Trait vertical à gauche de la proposition, de $V(I)$ jusqu'à $\mathfrak{S}_I/V(I)=\mathfrak{S}_{\widetilde I}$.]}

\margin{\uline{NB} $M_{0,I}^{!}$ se définit en \unsure{termes} de $I$, pour tt ens\unsure{.} $I$ de card\unsure{.} $\geq 3$, comme $M_{0,\unsure{n}}^{!}$}

\uline{Corollaire} Soit $i\in I$, alors l'action
de $\mathfrak{S}_{I\setminus\{i\}}\subset\mathfrak{S}_I$ sur $M_{0,I}^{!}$ est déduite
de celle de $\mathfrak{S}_{\widetilde I}$ via l'isom can $\mathfrak{S}_{I\setminus\{i\}}\xrightarrow{\sim}\mathfrak{S}_{\widetilde I}$

\scanpage{553}
\mspage{28\unsure{0|1}bis}
\deleted{\unsure{de}} des \unsure{s-groupes}
\[
 (31)\qquad \mathfrak{S}_3\times\mathfrak{S}_{n-3}\subset\mathfrak{S}_n
\]
de $\mathfrak{S}_n$ sur $M_{0,n}^{!}$, et pas de \unsure{celle|celles} de
$\mathfrak{S}_n$ tout entier. Dans le cas de $n=4$,
l'opération de $\Gamma'_2=\unsure{\mathrm{Sl}}(2,\mathbb{F}_2)\simeq\mathfrak{S}_3$ sur
$M_{0,4}^{!}$ qui nous intéresse, est celle qui
correspond au sous-groupe $\mathfrak{S}_{\unsure{1}}\times\mathfrak{S}_3\subset\mathfrak{S}_4$,
qui fixe le premier élément (\unsure{correspond.} à
la section nulle d'une courbe elliptique
--- \deleted{\unsure{qui}} et qui permute les trois \unsure{suivants}
\unsure{correspondts} aux trois éléments d'ordre $2$
d'une courbe ell. géom en \unsure{car.} $\ne 2$, qui
sont les trois autres \unsure{p\textsuperscript{t}} fixes de
$x\mapsto -x$, alors que l'opération \unsure{naïve|naive} de
$\mathfrak{S}_3$ sur $U_{0,3}$ correspond au ss-groupe
$\mathfrak{S}_3\times\mathfrak{S}_1$ de $\mathfrak{S}_4$, qui fixe les trois premiers
éléments. Petite question de gymnastique
\unsure{schématique} (\unsure{indépendante|indépendts} \unsure{directement|strictement} de la
géométrie des courbes elliptiques, mais
histoire de géom. projective de dim $1$) :
comment passer d'une de ces \deleted{\unsure{représentations}}\inserted{actions de}
$\mathfrak{S}_3$ à l'autre ? On a envie de penser que
ce sont les $\overline{\mathrm m}$s !

\scanpage{554}
Notons que si $I$ est un ens. à $4$
éléments, alors il lui est associé un
ens. à trois éléments des ``partitions de
type $(2,2)$'' $P_{2,2}(I)=\widetilde I$, et l'hom
\[
 (32)\qquad
 \diagram{\begin{tikzcd}[ampersand replacement=\&, row sep=small,column sep=large]
 \mathfrak{S}_I \arrow[r] \arrow[d,phantom,"\simeq"{rotate=90}]
   \& \mathfrak{S}_{\widetilde I} \arrow[d,phantom,"\simeq"{rotate=90}] \\
 \mathfrak{S}_4 \& \mathfrak{S}_3
 \end{tikzcd}}
\]
est surjectif --- son noyau est un sous-groupe
isomorphe (\unsure{non} can.\textsuperscript{t}) à $\mathbb{F}_2\times\mathbb{F}_2$ --- donc un
groupe commutatif \inserted{$V(I)$} annulé par $2$, i.e. un
vectoriel sur $\mathbb{F}_2$, et en \unsure{tt|tant} que tel de dim $2$.

En fait, $V(I)$ \unsure{opérant} sur $I$ fait de $I$
un $V(I)$ torseur --- donc $I$ est muni canoniquement
d'une structure \unsure{de|d'} \deleted{\unsure{espace}}\inserted{plan} affine sur $\mathbb{F}_2$,
et $V(I)\subset\mathfrak{S}_I$ est le groupe de ses
translations, $\mathfrak{S}_I$ \unsure{s'interprétant|s'interprète} comme le
groupe des autom. affines de $I$ (\uline{\unsure{tte}} permutation
de $I$ est affine --- il \deleted{\unsure{n'}} y a sur l'ens
$I$ $1$ et $1$ seule structure de plan affine sur
$\mathbb{F}_2$), et $\mathfrak{S}_{\widetilde I}$ comme le groupe $\mathrm{Aut}(V(I))$
--- ce qui suggère d'interpréter $\widetilde I$ comme
l'ens des trois éléments non nuls de $V(I)$, et
en effet, on constate que \deleted{\unsure{les}} l'on a \unsure{une|un}

\scanpage{555}
\mspage{282}

\unsure{provenant} de $I\setminus\{i\}\xrightarrow{\sim}\widetilde I$.

En particulier, pour $I=[1,2,3,4]$, pour
comparer l'action de $\mathfrak{S}_{\{1,2,3\}}$ et de
$\mathfrak{S}_{\{2,3,4\}}$ sur $M_{0,4}^{!}\simeq U_{0,3}$ --- où la première
est l'action \unsure{standard} de $\mathfrak{S}_3$ sur $U_{0,3}$,
la deuxième celle de $\mathfrak{S}_3$ sur $M_{1,1}[2]'$,
on utilise l'isom.
\[
\tag{38}
\mathfrak{S}_{\{1,2,3\}}\simeq\mathfrak{S}_{\{2,3,4\}}
\]
explicité dans (37), correspondant à la
bijection
\[
1\mapsto4,\quad 2\mapsto3,\quad 3\mapsto2
\]
des \unsure{ens.} d'indices, qui correspond donc
à l'\unsure{autom.}
\[
\tag{39}
\operatorname{int}(\sigma_1):\mathfrak{S}_3\longrightarrow\mathfrak{S}_3
\]
en identifiant maintenant les deux ens.
d'indices $[1,2,3]$, $[2,3,4]$ à $\{0,1,\infty\}$,
et en désignant par $\sigma_i\in\mathfrak{S}_{\{0,1,\infty\}}$
\inserted{(pour $i\in\{0,1,\infty\}$)} la
transposition des éléments $\ne i$. Donc

\uline{Corollaire} L'isom. canonique
\[
U_{0,3}\xrightarrow{\overset{\varphi}{\sim}}M_{1,1}[2]'
\qquad (\text{sur }\mathbb Z[1/2]),
\]
quand on fait opérer $\mathfrak{S}_3=\mathfrak{S}_{\{0,1,\infty\}}$ sur
l'un et l'autre membre, est compatible
avec \unsure{ses|ces} opérations, \uline{via} l'\unsure{autom.} $\operatorname{int}(\unsure{\sigma_1|\sigma_i})$, i.e.

\scanpage{556}

\[
\tag{40}
\varphi(u.x)=(\unsure{\sigma_1|\sigma_i}\,u\,\unsure{\sigma_1^{-1}|\sigma_i^{-1}})\varphi(x)
\]

Il faut quand $\overline{\mathrm m}$ démontrer la \unsure{prop. 4},
qui \unsure{équivaut} bien sûr à ceci

\uline{Corollaire} Soit $X$ une droite projective
relative sur \unsure{schém.} $S$, $I$ \unsure{ens} à quatre él.,
$(\deleted{\illegible{1}}s_i)_{i\in I}$ une famille de sections \unsure{mutuellement|mutuellemt}
disjointes, et $\sigma$ une \deleted{\unsure{transposition}}\inserted{\unsure{involution}} paire de
$I$ (correspondant à une partition $I=\{i_1,i_2\}\cup\{i_3,i_4\}$
par $\sigma i_1=i_2$, $\sigma i_2=i_1$, $\sigma i_3=i_4$, $\sigma i_4=i_3$). Alors $\exists$
\unsure{autom.} (évidemment unique) $u$ de $X$,
tel que $u\circ s_i=s_{\sigma i}$.

\uline{Dém} OPS $X=\mathbb P^1_S$, $\unsure{s_{i_1}|s_1}=0$, $\unsure{s_{i_2}|s_2}=\infty$, donc
$\unsure{s_{i_2}|s_2},\unsure{s_{i_3}|s_3}$ s'identifient à des sections de
$\underline{\mathcal O}_S^*$, on prendra \deleted{\illegible{1}} $u$ défini par
\[
u(z)=\lambda/\unsure{z|Z}\qquad\lambda\in\Gamma(S,\mathcal O_S^*)
\]
qui échange $0$ et $\infty$, et la condition
pour échanger $s_2,s_3$ sera $\lambda/s_3=s_4$, $\lambda/s_4=s_3$
i.e. $\lambda=s_3s_4$ (\uline{NB} il suffisait que
$(s_1,s_2,s_3)$ et $(s_1,\unsure{s_2|s_1},s_4)$ \unsure{sont|soient} mutuellement disjointes ---
pas la peine que $s_3,s_4$ soient mutuellement
disjointes $\ldots$)

\scanpage{557}
\mspage{283}

Plaçons-nous à nouveau sur $\mathbb Q$ (pour fixer
les idées), et considérons les groupes fondamentaux
des multiplicités modulaires $M_{1,1}$, $M_{1,1}'$.
On part de
\[
\tag{41}
\left\{
\begin{gathered}
\diagram{\begin{tikzcd}[ampersand replacement=\&, column sep=small,row sep=large]
M_{1,1}\arrow[d]\arrow[r,phantom,"\simeq" description]
  \& (M_{1,1}[n],\Gamma_n)\arrow[d]
  \& (n\geq3) \\
M_{1,1}'\arrow[r,phantom,"\simeq" description]
  \& (M_{1,1}[n],\underset{\substack{\Vert\\\Gamma_n/\pm1}}{\Gamma_n'})
  \& \begin{gathered}(n\geq2\unsure{)},\text{ en remplaçant}\\
      \text{pour }n=2\quad M_{1,1}[2]\text{ par}\\
      M_{1,1}[2]')\end{gathered}
\end{tikzcd}}
\end{gathered}
\right.
\]
qui \unsure{donne|donnent} des \unsure{isom.} de groupes extérieurs
\[
\tag{42}
\diagram{\begin{tikzcd}[ampersand replacement=\&, column sep=small,row sep=large]
\deleted{\illegible{1}}\pi_1(M_{1,1})\arrow[d]\arrow[r,phantom,"\simeq" description]
  \& \pi_1(M_{1,1}[n],\Gamma_n)\arrow[d]
  \& (n\geq3) \\
\pi_1(M_{1,1}')\arrow[r,phantom,"\simeq" description]
  \& \pi_1(M_{1,1}[n],\Gamma_n')
  \& \begin{gathered}(n\geq2\unsure{)},\text{ avec}\\
      \overline{\mathrm m}\ \text{\unsure{grain de sel}}\\
      \text{pour }n=2)\end{gathered}
\end{tikzcd}}
\]

comme $M_{1,1}[n]$ est connexe ($\overline{\mathrm m}$ s'il n'est
pas géom connexe), on trouve
\margin{c'est le th. de \unsure{Kronecker} d'irréductibilité \unsure{de l'équation|des équations} \unsure{cyclotomique|cyclotomiques}}
\[
\tag{43}
\left\{
\begin{gathered}
\diagram{\begin{tikzcd}[ampersand replacement=\&, column sep=small,row sep=large]
1\arrow[r,dash]
  \& \pi_1(M_{1,1}[n])\arrow[r]\arrow[d,"\unsure{\sim|2}"]
  \& \pi_1(M_{1,1}\deleted{\unsure{[n]}})\arrow[r,dash]\arrow[d]
  \& \Gamma_n\arrow[r]\arrow[d]
  \& 1 \\
1\arrow[r,dash]
  \& \pi_1(M_{1,1}[n])\arrow[r]
  \& \pi_1(M_{1,1}\deleted{\unsure{[n]}}')\arrow[r,dash]
  \& \Gamma_n'\arrow[r,dash]
  \& 1
\end{tikzcd}}
\end{gathered}
\right.
\]
qui permet de \unsure{réconstruire|reconstruire} $\pi_1(M_{1,1})$, en
tant qu'extension de $\Gamma_n$, quand on
\unsure{connait|connaît} $\pi_1(M_{1,1}')$ comme extension de
$\Gamma_n'$, par image inverse via $\Gamma_n\to\Gamma_n'$. On
trouve \deleted{\illegible{2}}\inserted{comme} de juste

\figure{[Trait vertical à gauche des deux lignes de la proposition.]}
\uline{Proposition} $\pi_1(M_{1,1})$ est une extension centrale
de $\pi_1(M_{1,1}')$ par $\{\pm1\}$.

D'autre part, on a les suites exactes,
\deleted{\illegible{5}}
d'homotopie sur $\operatorname{Spec}\mathbb Q$

\scanpage{558}

\diagram{(44)\quad
\begin{tikzcd}[ampersand replacement=\&, column sep=small]
1 \arrow[r,no head] \&
\pi_1(\overline{M}_{1,1})\inserted{{}=\unsure{\widehat{\mathcal T}_{1,1}^{+}}} \arrow[r] \arrow[d] \&
\pi_1(M_{1,1})\inserted{{}=\mathcal E_{1,1}=N_{1,1}} \arrow[r] \arrow[d] \&
\unsure{\Gamma_{\mathbb Q}} \arrow[r,no head] \arrow[d] \& 1 \\
1 \arrow[r,no head] \& \pi_1(\overline{M'_{1,1}}) \arrow[r,no head] \&
\pi_1(M'_{1,1}) \arrow[r,no head] \& \unsure{\Gamma_{\mathbb Q}} \arrow[r,no head] \& 1
\end{tikzcd}}

et le théorème transcendant fournit des
isomorphismes \inserted{ext.} canoniques

\diagram{(45)\quad
$\left\{
\begin{tikzcd}[ampersand replacement=\&, column sep=large,row sep=large]
\pi_1(\overline{M}_{1,1}) \arrow[d] \arrow[r,phantom,"\simeq" description]
 \& \mathrm{Sl}(2,\mathbb Z)^{\wedge} \arrow[d] \\
\pi_1(\overline{M'_{1,1}}) \arrow[r,phantom,"\simeq" description]
 \& \mathrm{Sl}(2,\mathbb Z)'^{\wedge}=\mathrm{Sl}(2,\mathbb Z)^{\wedge}/\mathord{\pm}1,
\end{tikzcd}\right.$}

\margin{Attention ne pas confondre $\mathrm{Sl}(2,\mathbb Z)^{\wedge}$ avec $\mathrm{Sl}(2,\widehat{\mathbb Z})$ !!!}

\deleted{\unsure{d'où} \illegible{1} fait que la suite}

On a \unsure{la} compatibilité essentielle
suivante entre (43) (44) (45) : la commutativité de

\diagram{(46) [Deux suites horizontales et trois flèches verticales descendantes.
La suite supérieure est
$1\longrightarrow\pi_1(\overline{M}_{1,1})\longrightarrow\pi_1(M_{1,1})\longrightarrow\unsure{\Gamma_{\mathbb Q}}\longrightarrow1$.
Sous $\pi_1(\overline{M}_{1,1})$, une identification verticale $\simeq$ avec $\mathrm{Sl}(2,\mathbb Z)^{\wedge}$, puis une flèche descendante portant «surj» vers $\mathrm{Sl}(2,\mathbb Z/n\mathbb Z)$.
La suite inférieure est
$1\longrightarrow\mathrm{Sl}(2,\mathbb Z/n\mathbb Z)\hookrightarrow\Gamma_n=\mathrm{Gl}(2,\mathbb Z/n\mathbb Z)\xrightarrow{\det}(\mathbb Z/n\mathbb Z)^*\longrightarrow1$.
Une flèche descend de $\pi_1(M_{1,1})$ vers $\Gamma_n$.
Une flèche descend de $\unsure{\Gamma_{\mathbb Q}}$ vers $(\mathbb Z/n\mathbb Z)^*$, portant «$\chi_n$ caractère cyclotomique».
À l'emplacement de $\unsure{\Gamma_{\mathbb Q}}$ et de la flèche qui le précède, une formule biffée :
\deleted{$\unsure{\Gamma_n}=\mathrm{Gl}(2,\unsure{\mathbb Z_n})$}.]}

Passant à la limite sur $n$, ceci donne
un diagramme commutatif

\diagram{(47)\quad
\begin{tikzcd}[ampersand replacement=\&, column sep=small]
1 \arrow[r,no head] \&
\bigl(\pi_1(\overline{M}_{1,1})\simeq\mathrm{Sl}(2,\mathbb Z)^{\wedge}\bigr) \arrow[r] \arrow[d,"\text{surj}"'] \&
\pi_1(M_{1,1}) \arrow[r] \arrow[d] \&
\unsure{\Gamma_{\mathbb Q}} \arrow[r,no head] \arrow[d,"\chi"] \& 1 \\
1 \arrow[r,no head] \& \mathrm{Sl}(2,\widehat{\mathbb Z}) \arrow[r,hook] \&
\mathrm{Gl}(2,\widehat{\mathbb Z}) \arrow[r] \&
\widehat{\mathbb Z}^{*} \arrow[r,no head] \& 1
\end{tikzcd}}

-- diagramme sur lequel nous allons revenir.

\scanpage{559}
\mspage{284}

Pour l'instant je vais expliciter le \inserted{dernier} isomorphisme (42), $n=2$ :
\[
(48)\qquad
\begin{aligned}
\pi_1(M_{1,1}')\inserted{{}=\deleted{\illegible{1}}N_{1,1}'}
&\simeq\pi_1(M_{1,1}[2]',\Gamma_2=\mathfrak S_3)\\
&\simeq\pi_1(U_{0,3},\mathfrak S_3),
\end{aligned}
\]
où le dernier isomorphisme correspond à
l'isom \inserted{$\operatorname{int}(\sigma_1):$} $\mathfrak S_3\simeq\mathfrak S_3$ explicité dans (40).

On trouve donc bien, comme prévu au
\S\ précédent -- et de façon entièrement
conceptuelle, l'isomorphisme \deleted{de suites}
\deleted{exactes} d'extensions

\deleted{(\unsure{49})\quad $\pi_1(M_{1,1}')\simeq\unsure{\widehat{\mathcal T}_{1,1}^{+\,\prime}}\simeq$}

\diagram{(49)\quad
\begin{tikzcd}[ampersand replacement=\&, column sep=small]
1 \arrow[r] \&
\pi_1(\overline{M_{1,1}'})\inserted{{}=\unsure{\widehat{\mathcal T}_{1,1}^{+\,\prime}}\simeq\mathrm{Sl}(2,\mathbb Z)'^{\wedge}} \arrow[r] \arrow[d,"\sim"'] \&
\pi_1(M_{1,1}')\inserted{{}\simeq N_{1,1}'} \arrow[r] \arrow[d,"\sim"'] \&
\unsure{\Gamma_{\mathbb Q}} \arrow[r,no head] \arrow[d,"\sim"'] \& 1 \\
1 \arrow[r] \& \pi_1(\overline U_{0,3};\mathfrak S_3) \arrow[r] \arrow[d,"\simeq" description,phantom] \&
\pi_1(U_{0,3},\mathfrak S_3) \arrow[r,no head] \arrow[d,"\simeq" description,phantom] \&
\unsure{\Gamma_{\mathbb Q}} \arrow[r,no head] \& 1 \\
\& \widehat\Pi_{0,3}' \& \mathcal E_{0,3}'
\end{tikzcd}}

D'ailleurs, repensant aux réflexions dans le
contexte transcendant, on voit que
cet isomorphisme

\diagram{(50)\quad
\begin{tikzcd}[ampersand replacement=\&, column sep=large]
\widehat\Pi_{0,3}' \arrow[r,"\sim"] \arrow[d,equal] \&
\pi_1(\overline{M_{1,1}'}) \arrow[d,"\simeq" description,phantom] \\
\widehat\Pi_{0,3}/\unsure{\{\ell_1^2,\ell_\infty^3\}} \&
\mathrm{Sl}(2,\mathbb Z)'^{\wedge}=\mathrm{Sl}(2,\mathbb Z)^{\wedge}/\mathord{\pm}1
\end{tikzcd}}

est associé à un isom
\diagram{(51)\quad
$\boxed{\begin{array}{ccc}
\Pi_{0,3}'&\xrightarrow{\sim}&\mathrm{Sl}(2,\mathbb Z)'\simeq\mathrm{Sl}(2,\mathbb Z)/\mathord{\pm}1\\
\Vert&&\\
\Pi_{0,3}/\{\ell_1^{\unsure{2}},\ell_\infty^3\}&&
\end{array}}$}

\scanpage{560}

déduit de l'isom de multiplicités analytiques
complexes
\[
(52)\qquad M_{1,1}^{\mathrm{an}\,\prime}\simeq
(U_{0,3}^{\inserted{\mathrm{an}}}\deleted{\illegible{1}},\mathfrak S_3)
\]

Il faudrait quand même expliciter
l'hom (51) -- qui n'est défini que
modulo autom. intérieurs, a priori,
par des formules explicites -- alors que
sa définition ici \unsure{est} de façon purement conceptuelle, ``géométrique'' (au sens
de la géométrie algébrique relative, des
courbes rationnelles et elliptiques sur des
espaces analytiques arbitraires). C'est là
un calcul-clef, \inserted{sûrement instructif,} qui ne devrait pas
présenter de difficulté particulière\ldots

\figure{[Court trait horizontal sous «difficulté», entre ce paragraphe et le suivant; fonction \unsure{de séparateur|de soulignement}.]}

J'ai un peu laissé tomber \inserted{en chemin} dans tout
ça le schéma modulaire ``grossier'' $\widetilde M_{1,1}$,
après avoir affirmé qu'il est $\simeq\mathbb E_{\mathbb Z}^{\unsure{1}}$, et
qu'il y a deux sections marquées,
dont les \deleted{points} val. en $\mathbb C$\unsure{,} correspondent
\unsure{aux} \inserted{deux} classes d'isom. de courbes elliptiques
\unsure{anharmoniques}. En termes de l'isom.
